\begin{prob}{290}{\hosi 8}{自作, 2023-7 大数宿題}
    $0\leq \theta < \frac{\pi}{2}$ を満たす実数$\theta$ が次の条件を満たすならば, $\theta = 0$ であることを示せ.

    $\tan{\frac{\theta}{2^n}}$ が有理数であるような非負整数 $n$ が無数に存在する.
\end{prob}

$\tan{2\theta} = \frac{2\tan{\theta}}{1-\tan^2{\theta}}$ より, $\tan{\theta}$ が有理数ならば $\tan{2\theta}$ も有理数である.
条件より, $\tan{\frac{\theta}{2^n}}$ が有理数であるような非負整数 $n$ が無数に存在するので, 上の議論からすべての $n$ に対して $\tan{\frac{\theta}{2^n}}$ が有理数である.
ここで, 互いに素な整数 $a_n, b_n$ を用いて $\tan{\frac{\theta}{2^n}} = \dfrac{b_n}{a_n}$ と置くとき,
半角公式より,
\[
\dfrac{b_n^2}{a_n^2} = \dfrac{a_n + b_n}{2a_n}
\]
である. 左辺は既約分数であり, 右辺の(自然数を用いた)分数を既約化すると左辺になることから
\[
a_n^2 \leq 2a_n
\]
がすべての$n$で成り立つ必要がある. $c_n = \dfrac{a_n}{2}$とおくと
\[
c_n \leq \sqrt{c_{n-1}} \leq \sqrt{\sqrt{c_{n-2}}} \leq \cdots \leq c_1^{1/2^n}
\]
が成り立つ. よって$a_n \leq 2c_1^{1/2^{n}}$である. 右辺は$N\to\infty$で$2$となるので,  $n$が十分大きければ$a_n \leq 2c_1^{1/2^n} < 3$が成り立つ. よって$a_n \leq 2$となる. \par
すべての$n\geq 0$に対して$0\leq \dfrac{\theta}{2^n} < \dfrac{\pi}{2}$であるので, $0<x_n<1$であることに注意する. すると, $a_n = 1$なら$b_n = 1$と決まる. つまり$x_n = 1$だから$\dfrac{\theta}{2^n} = 0$でなければならず, $\theta=0$を得る. \par
$a_n = 2$の場合, $b_n = 1$である($b_n = 2$の場合は互いに素であることに反する). このとき$x_n = \dfrac{1}{2} = \cos{\dfrac{\pi}{3}}$であるが, $x_{n+1} = \cos{\dfrac{\pi}{6}} = \dfrac{\sqrt{3}}{2}$は無理数であるので(1)に反する. \par
以上より$\theta=0$であることが示された. 実際にこの場合は, $\tan{\dfrac{0}{2^n}} = 0\in \Q$となって条件を満たす.
